\section{Weighted templates and simultaneous size readings}\label{sec:weighted}
A weighted template starts with disjoint finite permutation sets whose
cardinalities satisfy selected numerical relations.  Its defining subgroup
interval is itself finitely representable.  Our task is to recover those
relations from an arbitrary finite representation of a nontrivial lattice
image of the template.  The recovered cardinalities may differ from the
defining ones.  What must survive is their arithmetic, read simultaneously
from the parts of one partition.  The next section will bound the number
of lattice elements independently of the defining cardinalities.

We first pass to an ordinary subgroup interval, allowing either
orientation (\Cref{t:weighted-algebra}).  Normal and projection tests then
produce a forward interval in a nonabelian simple power, with nontrivial
base projections and product subgroups at its distinguished cut coatoms
(\Cref{t:weighted-product-success}).  Applying the forward marker theorem
to a pure-marker filter places the whole remaining template in an actual
alternating power (\Cref{t:weighted-restart}).  The cut intersections
recover a single partition of its permutation domain
(\Cref{t:weighted-alignment}); equal totals and squares are read from
that partition.  Only boundedly many markers are deleted, and the order
of each simple site remains at most the original ordinary top-group order.

Literal lattice calculations throughout use the original parameter order,
even when actual subgroup containment has the reverse order.  The
distinction between subgroups in the simple power and their ordinary
one-site counterparts is fixed in
\Cref{t:weighted-power,t:weighted-site-convention}.

\subsection{The template and its nontrivial images}

\begin{definition}[Units and stock]\label{t:weighted-units}
We use disjoint transitive permutation sets called units; their index set is denoted by \(T\), and
symmetric-group notation on a set of indices (when talking about permutations of letters) uses the
entire union of their units. There are singleton markers. On the units we take independent groups
with product \(Y\), of these kinds:
\begin{itemize}
\item full symmetric groups, including on the markers;
\item primary squares of side \(x\), with group \(A_x\times A_x\) slotwise, where \(8\mid x,\ x\ge16\);
\item auxiliary units of shape \((x/4)\times2\times2\) (same side-parameter restrictions), with \emph{full symmetric groups slotwise}.
\end{itemize}

For each primary side there is its auxiliary type (total size \(x\)), and for each axis length
used there are full reference units of that size. Units are supplied in cohorts of identical type
and size. Also, for each non-marker cohort so required (or basic extra full cohort used), include
a full offset cohort of size per unit one greater than its total size per unit. These extra
offsets themselves need no new offsets. In what follows every cohort, including markers, can be
required to have at least a sufficiently large number of members. For the group-structural tests
this lower bound can be \emph{absolute}, independent of sizes and number of cohorts. We explain
the finite losses involved along with the arguments. Assume there are nonmarker units in this use.
\end{definition}

\begin{definition}[Weighted template]\label{t:weighted-template}
Adjoin a foundation set \(Z\) larger than all the letters together. Use the ordinary interval
\[
       F=[S_Z\times Y,S_{Z\sqcup(\cup T)}].
\]
(Here \(\cup T\) denotes the union of the units.) The notation for its labels is
\[
       (I,U)=S_{Z\sqcup(\cup I)}\times U,\qquad
       I\subseteq T,\quad Y_{T\setminus I}\le U\le S_{T\setminus I}.
\]
The set \(I\) is the \emph{footprint}: its units have been absorbed into the foundation orbit.
\end{definition}

\begin{lemma}[Labels, cuts, and restrictions]\label{t:weighted-labels}
The labels in \Cref{t:weighted-template} are exactly the overgroups in \(F\).
Its coatoms are the labels \(C_D\) below.  Absorption in a join follows residual
orbits.  Axis filters and axis ideals again give the weighted templates on the
units that remain.
\end{lemma}
\begin{proof}
In any overgroup the conjugates of the supported \(S_Z\) generate the full symmetric group
supported on the orbit of \(Z\), by overlap (its size before expansion already exceeds half the
ambient degree). That orbit absorbs entire units. Absorption in joins thus follows orbits; in
particular the join of the \(x_j\)'s below will be top. We shall use the following labels:
\begin{align*}
 X_D&=(D,Y_{T\setminus D}), & x_j&=X_{\{j\}}, & M_D&=X_{T\setminus D},\\
 C_D&=(T\setminus D,S_D)\quad(D\ne\emptyset),& N&=C_T,
 & J^*&=(\emptyset,\prod_{j\in T}S_{\{j\}}).
\end{align*}
The \(C_D\)'s are precisely the coatoms (here and below test by absorption, or by the unequal
pieces of each cut with the large foundation). The filter above \(X_D\) and the ideal below
\(M_D\) have the templates on the remaining units (in the filter the foundation has grown); inside
\(M_D\) the independent unchanged \(Y_D\) factors out.
\end{proof}

\begin{convention}[Images and literal lifts]\label{t:weighted-images}
We will work even with an arbitrary nontrivial lattice image \(\bar F\) of \(F\) (surjective
lattice homomorphism); often the notation for labels refers also to their images. We call the ones
before the map literal labels. Intervals between comparable literal labels map onto the
corresponding intervals (join with the lower then meet with the upper to put a lift there); all
meets and joins may be tested by first doing the literal operation and then mapping.
\end{convention}

\begin{lemma}[Footprints and faithful pure links]\label{t:weighted-footprint}
In every nontrivial lattice image \(\bar F\), the footprint is well defined.
All cut coatoms survive distinctly and are all the coatoms of the image.
Residual orbit discrepancies are detected by absorption.  The pure marker
links, and canonical permutation intervals on copies of any fixed package of
units, are faithful.
\end{lemma}
\begin{proof}
The footprint \(I\) is unambiguous in the image. In fact a collapse with changed footprint gives a
comparable one (use meets), and for \(j\) gained it then identifies \(N\vee x_j=1\) with \(N\), by
meeting first with \(x_j\). Indeed \(x_j\) is under the upper label, while its intersection with
the lower is only residual (empty footprint). Thus the identification kills every \(x_R\) for full
units \(R\) by intersection with \(N\). For a grid unit \(O\), in each direction take a label
preserving a partition of \(O\cup R\) into the parallel lines of that direction plus a reference
block \(R\) of the matching size. Use full independent bases, \emph{alternating top} on the line
blocks including \(R\), and unchanged \(Y\) elsewhere, without absorption. This contains the
bottom: on parallel lines the grid groups induce even permutations (for the auxiliary grids both
other lengths are even). It joins with \(x_R\) to absorb \(O\). Intersecting all these labels with
\(x_O\) gives bottom: preserving each parallelism forces slotwise permutations and on the primary
squares the two alternating conditions enforce both prescribed slot groups. Thus all \(x_O\) are
killed as well, impossible. In particular the cut coatoms survive, distinctly with no new ones,
and orbit discrepancies of residual labels can be detected by joining with absorptions. Here
adjoining \(x_j\) to a label for an unabsorbed \(j\) absorbs exactly its residual orbit in
addition to the old footprint. On just markers, the resulting pure lattice, in particular in the
filter absorbing everything else, is exact: for a strict change of residual group (on singleton
positions) intersect down to a cyclic group generated by a discrepancy element, holding the
footprint fixed. A proper subgroup there has a smaller orbit on some cycle (an index divisor hits
the cycle length), detected by adjoining the absorption of one singleton. The same test works
inside a permutation label acting canonically on copies of a package of units modulo the
prescribed internal bottom actions.
\end{proof}

\begin{proposition}[From algebras to ordinary intervals]\label{t:weighted-algebra}
Every nilpotent join-preserving deflation of \(\bar F\) is zero.
If some nontrivial image of \(F\) is the congruence lattice of a finite algebra
of size \(m\), there is an ordinary finite-group interval representing the
dual of a nontrivial image of \(F\), whose top group has order at most
\(60^m m!\).  The full lattice \(F\) itself has a finite algebra representation.
\end{proposition}
\begin{proof}
Let \(f\) be a nilpotent join-preserving deflation.  It kills every atom:
otherwise it fixes that atom, contrary to nilpotence.  In particular,
\(f(x_R)=0\) for every full unit.  If \(A\) is any of the parallel-line
labels used in \Cref{t:weighted-footprint}, then \(x_O\le A\vee x_R\), so
\[
 f(x_O)\le f(A\vee x_R)=f(A)\le A.
\]
Also \(f(x_O)\le x_O\).  Intersecting over the directions therefore gives
\(f(x_O)=0\).  Since all the \(x_j\)'s join to the top, \(f=0\).
The idempotent-restriction and unary-nonunit argument of
\Cref{f:prop-ordinary-interval} applies starting with a least carrier representation
among all nontrivial images of \(F\). Idempotent restriction to more than one point would still
represent such an image since it gives a surjective lattice map; nonunits thus yield the nilpotent
join deflation as there. Consequently this minimal representation gives an ordinary finite group
interval for the \emph{dual} of a nontrivial \(\bar F\), by the diagonal construction in \Cref{f:eq-diagonal-interval}. Its top order is at most \(60^m m!\) if the representing size was \(m\), taking
the simple group \(A_5\). Also \(F\) itself is representable (the transitive coset action on its
ordinary interval).
\end{proof}

\subsection{Normals: absorption, parallelism, and signs}

We next show that a normal subgroup cannot produce an intermediate label
between the base and the top.  The tests successively exclude absorption
of units, movement between units, disruption of their parallel partitions,
and finally the remaining primary-square signs.  All calculations are
made with literal labels and then passed to the image.

\begin{convention}[Normal endpoint property]\label{t:weighted-normal-convention}
For an ordinary interval \([H,G]\), property \((*)\) means that every normal
subgroup of \(G\) is either contained in \(H\) or supplemented by \(H\),
and that modulo \(\operatorname{core}_G(H)\) there is a unique minimal normal
subgroup, a power of a nonabelian simple group, supplemented by the base.
In the next three results, \([H,G]\) represents \(\bar F\) in either orientation,
and \(e\) is the label of \(HL\), with \(L\trianglelefteq G\).
All inequalities between labels use the parameter order.
\end{convention}

\begin{lemma}[Projection and intersection tests for normals]\label{t:weighted-normal-test}
A normal-multiple label cannot satisfy
\[
 e\wedge b\le a<b\le e\vee a.
\]
Its axis projections are again normal-multiple labels.  If \(L\) is minimal
normal and nontrivial, a nontrivial proper intersection \(L\cap G_b\) cannot
be normalized at a family of labels whose actual groups generate \(G\).
The equations in the proof test precisely the normalizations needed below.
\end{lemma}
\begin{proof}
Write \(e\) for (the label of) a base times a normal, always reasoning with labels in the forward
\(\bar F\) order. Dedekind's law makes it left-modular: in particular \(e\wedge b\le a<b\le e\vee
a\) is impossible in either orientation. Projections to axis filters (by joining in the label
order) again come from normals: in actual order enlarging the bottom to \(K\) projects \(HL_1\) to
\(KL_1\), and shrinking the top to \(K\) projects it to \(HL_1\cap K=H(L_1\cap K)\), for
\(L_1\trianglelefteq G\). Axis intervals used after deleting a bounded number of units are still
nontrivial images on the corresponding templates, by interval surjectivity and footprint
detection. For the normal tests below we suppose that \(e\) is a proper nonendpoint.

Here is also the intersection test used repeatedly below in the normal argument. Write \(G_b\) for
the actual ordinary overgroup at label \(b\), and \(G_e=HL\) for the actual normal multiple,
\(L\trianglelefteq G\) (similarly after any indicated base-core passage). If \(H<G_b<HL\), then
\(I_b=L\cap G_b\) is nontrivial proper in \(L\), by \(G_b=H I_b\). It is normalized at every \(s\)
for which \(\langle G_b,G_s\rangle\cap HL=G_b\), since then it is just \(L\cap\langle
G_b,G_s\rangle\). The sufficient equation in parameter order is thus \(e\wedge (b\vee s)=b\)
forward, \(e\vee(b\wedge s)=b\) reverse. If \(L\) was taken minimal normal nontrivial and the
indicated \(G_s\)'s generate \(G\) (label join forward, label meet reverse), this normalization is
impossible.
\end{proof}

\begin{lemma}[A proper normal label is internal]\label{t:weighted-normal-internal}
If \(e\) is a proper nonendpoint normal-multiple label, then
\(I(e)=\emptyset\) and \(e\le J^*\).
The conclusion remains available after any of the bounded axis restrictions
used in the proof, and when all unit actions are replaced by full symmetric
groups.
\end{lemma}
\begin{proof}
\noindent\textbf{Absorption and the offset comparisons.}
Take four disjoint equal-size unions of units in a cyclic order.
Absorb the outside and allow independent full symmetric groups on
the four blocks.  Adjoining the square of their rotation gives a
label \(a\), and adjoining the whole rotation gives \(b>a\), with
strictness detected by orbits.  Suppose the four counts of absorbed
letters from \(I(e)\) are nonconstant, with a positive count in each
parity class of positions.  An odd rotation cannot preserve these
counts, so \(e\wedge b\le a\).  But the square rotation connects
each parity pair, both of which meet the absorbed part of \(e\);
hence \(e\vee a=1\).  This is forbidden.

In any cohort, therefore, either every unit is absorbed or at most
one is: otherwise put two absorbed copies in opposite parity
positions and a nonabsorbed copy in one of the remaining positions.
Each sufficiently numerous cohort consequently has a majority
absorption bit.  Compare a typical size-\(s\) unit together with
a typical marker against a typical size-\(s+1\) offset unit.
Write their majority bits as \(\epsilon,\mu,\eta\), respectively.
The four-block test, using both block types and repeating a positive
one in both parity classes if their counts differ, forces
\[
 s\epsilon+\mu=(s+1)\eta.
\]
For bits in \(\{0,1\}\), this implies
\(\epsilon=\mu=\eta\).  Every required cohort and its offset
therefore have the marker cohort's bit.  If that bit were one,
the preceding cohort dichotomy would make every unit absorbed,
contrary to \(e<1\).  All majority bits are zero.

Finally, two absorbed units from different cohorts are impossible.
Use four blocks each containing one copy from each of those cohorts,
and put the two absorbed copies in different parity classes, choosing
the other copies nonabsorbed.  The same test applies.  Thus
\(|I(e)|\le1\).  These arguments use only sizes, offsets and
boundedly many copies, so they also hold with full unit actions
and after the indicated bounded axis deletions.

\noindent\textbf{Excluding a singleton footprint.}
If \(I=\{O\}\) and the residual group moves between other units, shift above \(x_j\) at the source
of such a move. At least two absorptions remain (\(O\) and another unit hit from \(j\)), so the
shift is top by the count just proved (the same count in a filter deleting a unit), meaning the
residual group is transitive on the complement. On four copies including \(O\) the square test
then gives total absorption on join anyway (we reach another unit besides \(O\), then its whole
transitive complement orbit), impossible. If instead the residual group preserves units
separately, use three copies including \(O\) with rotation group \(A_3<S_3\) on the blocks,
outside absorbed. This increase is strict (meet with the version allowing just one transposition
on blocks, then detect orbits); intersection with \(e\) cannot move the blocks while join absorbs
them. So \(I=\emptyset\).

\noindent\textbf{Excluding residual mixing.}
By shifting once at a mixed unit and applying the same conclusions,
mixing units now forces transitivity on all the residual letters (there is an extra absorption in
that shifted label, hence it is the filter's top). These nested temporary shifts delete only
boundedly many units before doing the counts.

For such a transitive \(e\), every marker \(x_i\) complements it. Normal complementation gives
isomorphisms in actual subgroup order (when \(K\) complements \(HL\), \(L\trianglelefteq G\),
project \(K\) to \(G/L\) to identify \([H,K]\) with \([HL,G]\)). Forward this forces \(e=N\) by
marker atomicity (it makes \(e\) a residual coatom). Replace the base by the label \(J^*\), i.e.
switch to all full unit groups, with the same transitive projected label. In the new base-core
quotient take a minimal normal inside the normal giving that label; it is still missing and
proper, so still transitive by what was proved thus far (nonmixing would now mean bottom), hence
still maps to \(N\). For a residual marker transposition \(d\) over the new base, we have
\(0<d<N,\ N\wedge(d\vee x_j)=d\) for every unit, using notation relative to the full-unit base.
Thus the intersection of the minimal normal with the subgroup labeled by \(d\) is nontrivial
proper and normalized by all labels \(x_j\) which generate the top, impossible. Indeed inside an
actual base times that normal, all overgroups of the base factor as the base times their own
intersections with the normal; and here enlarging \(d\) by any \(x_j\) leaves its intersection
unchanged.

In reverse orientation the complementation identifies \([x_i,1]\) with \([0,e]\) by meeting with
\(e\) in \(\bar F\). Thus \(X_{\{i,k\}}\wedge e>0\) on distinct markers, forcing every marker
transposition label beneath \(e\) (the literal residual action in that intersection only has the
two singleton markers available over bottom). A supported marker three-cycle group including
\(i\), which lies below \(e\) by these transpositions, then cannot be in the indicated meet image:
joining it with \(x_i\) absorbs all three and they admit transpositions below \(e\). These small
permutation labels over bottom are exact by orbit detection. So in either orientation every proper
nonendpoint normal label lies under \(J^*\).
\end{proof}

\begin{lemma}[Parallel restrictions on an internal normal]\label{t:weighted-normal-parallel}
After quotienting by the base core, suppose that a minimal nontrivial normal
subgroup gives a proper nonendpoint label \(e\).  Let \(B\) be the intersection
of \(J^*\) with all full symmetric parallel-block labels.  Then \(e\le B\).
This case is impossible in reverse orientation.  In forward orientation
\([0,B]\) is the exact lattice of subspaces of the sign space having two
coordinates for each primary square.
\end{lemma}
\begin{proof}
Pass now to the original base-core quotient and, if a proper nonendpoint occurs, take the normal
to be minimal nontrivial therein (within the purported normal), written \(L\). This still gives a
proper nonendpoint internal \(e\) (actual base times \(L\)); internal representatives can be used
since \(e\le J^*\).

For a parallel direction on \(O\), write \(A<W\) for the alternating-top parallel label used
earlier and its full symmetric-top version, now allowing full internal groups on other units. This
increase survives: adjoin a transposition \(d\) on two spare markers to both, then intersect with
a graph label coupling this swap to odd line-block action, giving different marker orbits,
detectable by absorption. Suppose \(e\not\le W\). Some element of an internal representative then
disrupts the partition; conjugating supported full bases and using primitive alternating top
generates independent full action there (a cross-transposition in addition to the within-block
symmetric groups suffices), so \(e\vee A\ge W\). Thus \(D=e\wedge W\not\le A\) by left-modularity.
Forward, \(0<D<e\) in this situation, and
\[
                  e\wedge(x_k\vee D)=D\qquad(k\ne O),
\]
since \(x_k\), like the literal internal intersection giving \(D\), preserves \(O\) and its
indicated lines; on meeting back with the internal representative one thus stays under \(W\).
Hence the nontrivial proper intersection of \(L\) at \(D\) is normalized both at \(W\) (where the
intersection is the same) and all these \(x_k\) (use \(x_k\vee D\)); together they absorb
everything, impossible by minimal normality.

In reverse orientation in fact no internal element of a representative for \(e\) can preserve the
direction's lines on \(O\) and act oddly on them. Otherwise set \(K=e\vee d\) with \(d\) just a
spare marker swap over bottom. This is strictly above \(e\) (orbit test) but residual. Take the
graph label \(P'\) coupling this swap to odd line-top on \(O\), otherwise allowing internal
actions preserving units and this parallelism. We have \(d\wedge P'=0\), and \(e\vee(K\wedge
P')=K\) using the presumed odd element times the swap. Thus the intersection of \(L\) at \(K\),
nontrivial proper by \(e<K<1\), is normalized at \(d\) (by \(K=e\vee d\)) and at \(P'\) (enlarging
actually, i.e. using \(K\wedge P'\), leaves the intersection unchanged). These generate actual top
by \(d\wedge P'=0\). This too is impossible; it rules out in particular \(D\not\le A\) obtained
above under a disruption (an odd top element in the internal intersection with \(W\) fixes the
reference block, hence acts oddly on the lines of \(O\)).

It follows in both cases that \(e\le B\), the intersection of \(J^*\) with all the \(W\)'s. This
is slotwise full on the grids and differs from bottom only by primary square sign bits. Reverse
has now been excluded (meet to place the representative under literal \(B\), then use the odd-line
test).

In the forward case \(0<e\le B\), so the image of the literal sign-space
interval is nontrivial.  The subspace lattice is simple: a nontrivial
congruence collapses an atom by atomisticity, and every two distinct atoms
lie in a plane, where such a collapse propagates to all atoms.  Hence this
interval is exact.
\end{proof}

\begin{proposition}[Normal endpoint property for weighted images]\label{t:weighted-normal}
For sufficiently large absolute cohort stock, every ordinary interval
representing a nontrivial weighted image, in either orientation, satisfies
\((*)\).  Its base acts transitively on the simple factors of its monolith
modulo the base core.  The same assertion holds for the bounded axis deletions
used here, provided the remaining stock is sufficient.
\end{proposition}
\begin{proof}

By \Cref{t:weighted-normal-internal,t:weighted-normal-parallel}, only the
forward sign-space case remains.  Take a minimal missing normal inside any
purported proper normal multiple after passing to the base-core quotient.

Forward, \([0,B]\) must now be an exact subspace lattice over \(\mathbf F_2\), since that literal
interval is simple and its image is nontrivial (\(0<e\le B\)). Here the slotwise groups over
literal bottom have quotient given by the sign space with two bits per primary square. Subspace
lattices here are simple, e.g. by atomisticity and propagation of any atom collapse across the
planes. For an active bit of the sign space of \(e\), meeting with its parallel alternating-top
test \(A\) gives its kernel \(D'\). If this is nonzero, the same formula \(e\wedge(x_k\vee
D')=D'\), \(k\ne O\), and normalization at \(A\) give the earlier forward contradiction. Thus
\(e\) has sign dimension one with vector \(v\ne0\). Swapping the two coordinates of any primary
square by \(\tau\) must fix \(v\): otherwise use as \(a\) the label adjoining \(\tau\) to bottom
and as \(b\) this label extended by the sign line \(\langle v+\tau v\rangle\); \(e\wedge b\le
a<b\le e\vee a\), strictness tested under \(B\). Hence at an active square both bits occur
together.

Choose an active square and use its \(x\) lines in one direction,
together with its size-\(x\) auxiliary unit as one additional block.
Let \(a\) have independent alternating groups inside these blocks and
the alternating group on their \(x+1\) positions.  Match coordinates
along the lines and choose any matching with the auxiliary block.
Outside this union allow full internal groups, without absorption.
The prescribed bottom lies in \(a\).  Let \(b\) instead allow all base
sign vectors of even total, with the same alternating top.  Then
\(a<b\): the base sign which is odd on every original line and even
on the extra block lies in \(b\), and its nontriviality is detected by
meeting with the exact sign interval \([0,B]\).

Write \(\alpha\) for the line-position sign and \(\beta\) for the
repeated within-line sign of the chosen square.  The preceding swap
test gives \(\alpha(v)=\beta(v)=1\).  An element of \(e\wedge b\)
has \(\alpha=0\), since \(b\) has alternating top.  As \(e\) is the
single sign line \(\langle v\rangle\), its sign component must
therefore vanish, including \(\beta\).  The auxiliary bottom action
is also even on its letters: its first-slot permutations repeat four
times, and each binary-slot swap repeats an even number of times.
Consequently \(e\wedge b\le a\).

On the other hand, a representative of the nonzero sign of \(e\)
has base-sign vector
\[
 u=(1,\ldots,1,0)
\]
on these \(x+1\) blocks, and may also permute the original lines.
Conjugate the supported matched alternating top by this representative,
then remove the top permutations using \(a\).  Modulo the alternating
bases this produces \(u-hu\) for every even block permutation \(h\).
These differences span the augmentation space: moving the unique zero
gives every pair of coordinate positions.  The conjugations remain
supported on this union regardless of the representative's actions
on other units.  Hence \(e\vee a\ge b\), contradicting the
normal-multiple test.  This proves the normal endpoint property.

The lattice is nonmodular (the exact pure marker link includes the subgroup lattice of \(S_4\),
already nonmodular: with a 4-cycle \(r\) and a dihedral reflection \(s\) inverting it, \(\langle r
s\rangle\) has trivial meet with \(\langle s,r^2\rangle\) but joins with \(\langle s\rangle\)
above it). For completeness, in the faithful base-coset action every minimal normal is
transitive by the endpoint property.  An abelian minimal normal would identify
the interval with an invariant-subgroup interval in an abelian group, which
is modular.  Two distinct minimal normals commute; supplementation by one
then makes the base induce all inner actions on the other.  Intersections
in the latter are invariant normal subproducts, again giving a modular
interval.  Thus there is a unique nonabelian monolith modulo the base core,
and it is supplemented.  The base is transitive on its simple factors,
because the full group is and the monolith acts internally on each factor.

The stock bounds thus far only spent boundedly many copies for parallel tests, rotations, and
nested shifts deleting at most a bounded number of units.
\end{proof}

\subsection{Product closures in both orientations}

\begin{convention}[Actual groups and product closure]\label{t:weighted-power}
After passing to the base-core quotient we thus have \(G=H\mathcal W\), \(\mathcal W=S^s\) a
minimal normal nonabelian power, and the labels are exact \(H\)-invariant subgroups \(K_b\) in
\([P,\mathcal W]\), where \(P=H\cap\mathcal W\), transitive operator action on the sites. Thus
\(K_b\) is written on the simple-power side (actual subgroups over the interval base correspond by
intersection with \(\mathcal W\)), not a formal permutation label. In reverse, label joins mean
actual intersections, and label bottom means actual \(\mathcal W\). We use \(c\) for coordinate
product closure, and call these its product labels when closed. Recall the extension description
in \Cref{f:lem-extension-description} of subdirect labels (site stabilizer \(D_0\le H\), action image \(\Lambda\le
{\rm Aut}(S)\), extensions to overgroups of \(D_0\) in \(H\)); above a given subdirect group in
actual order all are just restrictions on its domain interval, by reverse correspondence. Also,
intersections of comparable distinct subdirect actual groups with a fixed invariant product must
differ if nontrivial: the restrictions are independent across the tie parts, with nontrivial
projections at every site by operator transitivity, so a strict tie change distinguishes them. And
a subdirect joined (in actual order) with a nontrivial invariant product gives the whole power, by
conjugating independently supported elements and simplicity. These facts and the projection
arguments below can use the same transitive operator action on restricting a filter of actual
groups (or add the enlarged base to the operators by inner actions).
\end{convention}

\begin{proposition}[Forward proper projections]\label{t:weighted-forward-projections}
In a forward realization \([P,\mathcal W]^H\), suppose \(P>1\) and
\(c(P)<\mathcal W\).  Every cut coatom is a product label.  If
\(P<c(P)\), its closure label has empty footprint and lies below \(J^*\).
\end{proposition}
\begin{proof}
Forward, suppose first \(P<c(P)<\mathcal W\), and put \(K_p=c(P)\). Joining a label with \(p\) to
top forces subdirectness, which also suffices since \(K_p\) contains nontrivial independent
support. Thus a test \(p\wedge b\le a<b\), \(p\vee a=1\), is impossible by strictness on
intersections (both \(K_a,K_b\) then subdirect, intersections containing \(P>1\)). Exactly the
same four-block rotation test gives \(|I(p)|\le1\). In the singleton case \(I=\{O\}\), if any
other copy of \(O\) is moved in the residual action to touch a different unit, place it opposite
in parity from \(O\) in the four-copy test, with destination outside the chosen four. The join
again absorbs everything. If none are thus moved use the three-copy test instead. Hence
\(I=\emptyset\).

Shift above \(x_j\); if its closure is not full the closure label has empty extra footprint there
by the same tests (or trivially when \(x_j\) is already product). So if \(p\) mixes \(j\) with
other units then \(x_j\) must project fully by monotonicity of closure, and \(x_j\vee p=1\),
forcing transitivity of the residual group. In that case all marker \(x_i\)'s are subdirect and
their intersections with \(p\) are bottom. Meeting also with \(X_{\{i,k\}}>x_i\) forces all marker
transposition labels beneath \(p\) by strictness. The \(K_{x_i}\)'s have identical ties and
identical twists restricted to the projections of \(P\), by the common intersection \(P\) with
\(K_p\). Therefore the images of all the filters above \(x_i\) under meeting with \(p\) must agree
(ties may be split by any invariant refinement, retaining the twists, then taking the product of
the correspondingly restricted diagonals on intersection with \(K_p\)). Yet a supported marker
three-cycle group \(d'\) lies beneath \(p\) by the transpositions, and occurs in the meet image
for \(i\) outside its three positions, as \(p\wedge(x_i\vee d')\) (\(p\) is residual, \(i\) is
singleton). It cannot occur for \(i\) inside by the same absorption contradiction as before. Thus
\(p\le J^*\). Now every cut coatom is product, since subdirectness there would join to top with
\(p\). The same cut conclusion holds immediately if \(P=c(P)\ne1\).

To elaborate on the common meet-image test in that argument, put \(\mathcal T=c(P)=\prod T_o\)
with the index here ranging over sites. A subdirect \(K_{x_i}\) meeting \(\mathcal T\) in \(P\)
restricts on each tie part to the full twisted diagonal of the corresponding \(T_o\)'s in that
intersection: the projection there is exactly \(T_o\) since it comes from \(P\). These are
independent across parts. Thus whether two sites are tied, and if so the restriction of their
identifying isomorphism to \(T_o\), can be read from the two-site projection of \(P\) (graph of an
isomorphism versus a whole product; \(T_o\ne1\)). Above \(K_{x_i}\) one may split by every
operator-invariant refinement, preserving its twists within the remaining ties; this is again
invariant by equivariance of the original ties, and these are all the invariant overgroups in the
power here. Their intersections with \(\mathcal T\) depend only on that same restricted-twist data
and the refinement. This proves the claimed agreement even if the full twists were not specified
by \(P\). In marker notation with \(i\) outside the triple one can bound \(p\wedge(x_i\vee d')\)
also by \(N\wedge(x_i\vee d')=d'\) calculated literally, and it contains \(d'\) in the image.
\end{proof}

\begin{definition}[Product success]\label{t:weighted-success-definition}
We say that \emph{product success} holds when a forward realization in a
nonabelian simple power has nontrivial projections of its base intersection
at every site and every cut coatom is a product subgroup.  We allow passage
to this configuration by deleting one marker in an axis filter.
\end{definition}

We next establish in \emph{reverse} orientation that for \(P>1\) we must have \(c(P)=\mathcal W\),
and for \(P=1\) all nonbottom actual labels are subdirect. The \(P>1\) argument only needs an
invariant interval with those base and action conditions.

\begin{lemma}[Subdirect copy actions in reverse orientation]\label{t:weighted-copy-action}
In a reverse invariant realization with \(P>1\), let \(z\) canonically
permute \(k\) identical nonempty packages of units, acting as the prescribed
bottom inside each package.  For a sufficiently large absolute \(k\),
\(K_z\) is subdirect.  The same holds for its canonical alternating-copy
label once that label is sufficiently large.
\end{lemma}
\begin{proof}
Take disjoint identical nonempty packages of units and a residual label \(z=Y\rtimes S_k\)
permuting \(k\) copies canonically (fix coordinated matchings of corresponding units respecting
their types and actions). Foundation actions are implicit in residual notation; by a canonical
subgroup label here we mean the inverse image in \(z\) of a subgroup of \(S_k\). The interval
\([0,z]\) in label order is exact \({\rm Sub}(S_k)\), by the cyclic discrepancy and orbit test on
these copies. We claim \(K_z\) subdirect for sufficiently large absolute \(k\), here when \(P>1\).
Otherwise its closure corresponds to a nonidentity subgroup \(U\le S_k\). Labels indexed by
subgroups of \(U\) generated by prime order elements are product, since the corresponding atoms
(actual coatoms above the nontrivial product \(K_U\)) must be so. Hence \(U\) cannot contain a
large supported \(A_v\): that subgroup is prime-generated and would give a product base
\(K_{A_v}\ne1\) for an exact reverse \({\rm Sub}(A_v)\) inside this invariant interval, contrary
to the reverse alternating exclusion of \Cref{r:reverse-exclusion}.

If \(a<c_0\) in \({\rm Sub}(S_k)\) and \(c_0\cap U=1\), their labels are subdirect (by
monotonicity, the label of each closure is below \(U\) and below \(c_0\)), so \(\langle
a,U\rangle<\langle c_0,U\rangle\) by intersection strictness with \(K_U\). Thus any transposition
moved by conjugation by \(u\in U\), together with its conjugate, generates a group meeting \(U\)
nontrivially. We get some prime-order \(h\in U\) supported on at most four letters. Match a
nontrivial orbit of \(h\) to a fresh ordered set of the same size (2 or 3) and exchange them
simultaneously pairwise by an involution. By the very same strict-join test its group together
with its \(h\)-conjugate (dihedral, two such exchanges) meets \(U\) nontrivially. Any nonidentity
element there rotates the two sets nontrivially or exchanges them in matched fashion. By finite
Ramsey homogenize a large reservoir of fresh points so a constant such pattern on increasing
tuples lies in \(U\). Comparing two tuples differing at just the last entry \(y,y'\) gives a
supported three-cycle by taking quotients: the two permutations differ by conjugation by
\((y\,y')\) and the first (say \(w\)) fixes \(y'\) and sends \(y\) to a common anchor (possibly of
the original orbit) distinct from \(y\). Thus one quotient is \((w(y)\,y')(y\,y')\). This anchor
is fixed as a choice of position when the prefix is fixed and the last entry varies. Varying these
cycles gives a large supported alternating group, contradiction. All these Ramsey requirements are
absolute (orbit lengths two or three, bounded many permutation patterns, with sufficiently many
points for the absolute reverse alternating exclusion). This proves the symmetric-copy claim.  Its
alternating-copy subgroup also
has a subdirect actual group: in reverse order that group contains the
subdirect group just obtained.
\end{proof}

\begin{lemma}[Small labels are subdirect in the reverse proper-closure case]\label{t:weighted-small-labels}
Suppose \(P>1\) in reverse orientation and
\(K_p=c(P)<\mathcal W\).  Then \(I(p)=\emptyset\), and every label supported
on any fixed bounded number of units is subdirect.  For the subsequent tests,
a support bound of ten is enough.  The required cohort stock is absolute.
\end{lemma}
\begin{proof}
Suppose now \(K_p=c(P)<\mathcal W\). We have \(I(p)=\emptyset\): for any absorbed \(O\) use such a
canonical permutation label on many copies including it; its alternating and symmetric versions
are distinct subdirect labels, yet have the same joins with \(p\) (all copies are absorbed),
impossible.

Call a label small here when supported on a bounded number of units (literally below their
\(X\)-label); we need only fixed bounds, say up to ten. We show that every such label has a
subdirect actual group. Otherwise an atom \(e_i\) (actual coatom) beneath its nonzero closure
label in parameter order must be product, and \(e_i\le p\): the actual coatom contains the
nontrivial product closure, so cannot itself be subdirect, hence is closed by maximality. It is
residual and lifts to a literal label supported on the package in question at empty footprint.
Take there a literal cover \(B_0<C_0\) of groups giving the images bottom and \(e_i\) (groups
written on just the package). Choose many disjoint clones and write \(e_j\) for the analogous
images of \(C_0\); they too are distinct atoms and the lowers collapse. Indeed image congruences
within residual package intervals transfer both ways by canonical swaps: join with the swap label
over \(Y\), then meet with the full residual symmetric label on the other package (only \(Y\)
elsewhere). Literally this copies a supported group from one package to the other, by independent
support and conjugation. Thus both collapses and noncollapses there transfer, and each supported
interval maps onto its image interval. Distinctness of \(e_j\)'s follows also by their bottom
meets.

\noindent\textbf{Transporting product atoms.} Every \(e_j\) must likewise be product. Indeed a coatom if not product is subdirect. If subdirect here for \(j\ne i\), it contains (in actual order) both intersections of \(K_{e_i}\) with the actual groups at the canonical labels of \(\langle(i\,j)\rangle\) and \(\langle(i\,j\,l)\rangle\) on copies. This is because on the parameter side, joining \(e_i\) with either permutation label contains \(e_j\) by conjugating the literal \(C_0\). Thus its ties refine both sets of permutation-label ties, since the indicated intersections with the product \(K_{e_i}\) project nontrivially and independently across each set of parts. But their common refinement is discrete: the two canonical labels are subdirect with a common subdirect lower group \(K_z\), hence consistent twists, and generate actual top (the canonical parameter subgroups meet trivially); any sites tied in both would keep that common twist on joining. This contradicts proper subdirectness of \(K_{e_j}\). Being product and above \(P\) as an actual subgroup, each is above \(c(P)\), i.e. \(e_j\le p\).

\noindent\textbf{The faithful product interval.} Use \(k\) large such packages. Between the products of the literal covers the interval is Boolean if \(B_0\) is not normal in \(C_0\) (any full projection conjugates the supported \(B_0\) to generate \(C_0\) there); otherwise a subspace lattice over \(\mathbf F_\ell\) for the prime order of \(C_0/B_0\). This interval maps faithfully to \([0,\bigvee e_j]\) by survival or simplicity, and all entries are product (the join of the \(e_j\)'s lies under \(p\), hence all atoms under it too must be coatom products in actual order; now use atomisticity). Take \(k\) coprime to \(\ell\) in the latter case (one of two consecutive large counts suffices). Let \(E\) be the augmentation hyperplane in that case, the whole join in the Boolean case. It is nonzero and no atom beneath it is fixed by the canonical alternating copy action (label \(a\) over \(Y\)), using the faithful literal product interval. For the hyperplane an invariant line would have a fixed nonzero vector by perfection of \(A_k\), hence constant on the copies, impossible by \(\ell\nmid k\).

\noindent\textbf{The extension image.} For the subdirect \(K_a\), write \(D_a\) for its extension domain with map \(\beta\) to \({\rm Aut}(S)\). By restrictions \([D_0,D_a]\) realizes exactly \({\rm Sub}(A_k)\), so the kernel of outer action must be supplemented by \(D_0\) (the normal test in \Cref{sec:reverse}; if retained, the base-core quotient would be soluble by Schreier, impossible there). Thus \(\beta(D_a)=(\beta(D_a)\cap S)\Lambda\), identifying \(S\) with inners. Also
\[
                   (a\vee E)\wedge p=E
\]
by factoring the literal join as a set product (the lift of \(E\) is normalized by the copy
action) and using a representative of \(p\) above that lift, since \(a\wedge p=0\). Here \(a\wedge
p=0\) comes from joining a subdirect with a nontrivial product in actual order, and literally the
set product intersected with that representative of \(p\) is the lift of \(E\) times the
intersection with the lift of \(a\). Hence the product closure of \(K_a\cap K_E\) is \(K_E\) (the
equality gives generation with \(c(P)\), already under the closure by \(P\le K_a\cap K_E\)).
Writing \(T_e={\rm proj}(K_e)\) at the site, this means \(\beta(D_a)\) normalizes \(T_E\). Indeed
intersection with a product has full projection onto it precisely when its factors are matched by
the diagonal twists; by transport along the tie through the site this is precisely invariance of
the site's factor under the extension image.

Here \(T_E\) is nontrivial and proper, by \(P>1\), site transitivity
and the proper product \(K_E<\mathcal W\).  Its normalizer in \(S\)
is proper by simplicity, and is invariant under \(\Lambda\).  Transporting
it gives a product label \(0<f\le E\).  Choose an atom \(h\le f\),
and put \(B_a=\beta(D_a)\cap S\).  The normalization just proved and
the reversal of label order give
\[
 B_a\le N_S(T_E)=T_f\le T_h.
\]
Also \(\Lambda\) normalizes \(T_h\), as it does every invariant
product projection.  Thus
\(\beta(D_a)=B_a\Lambda\) normalizes \(T_h\).  The same
intersection criterion now gives
\[
 c(K_a\cap K_h)=K_h.
\]
The alternating copy action moves the literal lift of \(h\) to a
different atom \(h'\) under \(E\), so \(h'\le a\vee h\) in label
order.  Its actual group is a product and contains \(K_a\cap K_h\);
it must therefore contain the displayed product closure \(K_h\).
This contradicts the distinct actual coatoms \(K_h,K_{h'}\), proving
the small-label claim.

In this extension-image argument, \([D_0,D_a]\) really is the whole ordinary interval: all actual
overgroups of the subdirect in the invariant interval are still subdirect and correspond exactly
to restrictions of its extension map. For the normalization criterion at the tie through a site
\(o\), at \(t o\), \(t\in D_a\), the product factor of \(K_E\) is the transport by \(t\) of
\(T_E\), whereas the entry matched by \(K_a\) to \(u\) at \(o\) is the transport of
\(\beta(t)^{-1}(u)\). Thus all \(u\in T_E\) occur in the intersection exactly for stability under
\(\beta(D_a)\), and similarly for the product at \(h\); other tie parts transport. Finally the
``different atom'' test only conjugates a literal lift within the displayed Boolean/subspace
interval, where alternating copy permutations preserve the lift of \(E\) and move the given atom's
lift to a different atom there. The distinctness survives by faithfulness on this interval, and
the conjugate is contained in the literal join with the copy label, so the conclusion in the image
does not assume invariance of the entire quotient map under copy permutations.
\end{proof}

\begin{lemma}[Residual products are impossible in reverse orientation]\label{t:weighted-residual-obstruction}
Suppose all labels on the bounded supports used below are subdirect in a
reverse realization.  There is no nontrivial proper product actual group
whose parameter label \(p\) has empty footprint.  This assertion applies
both when \(P>1\) and when \(P=1\).
\end{lemma}
\begin{proof}
For small residual \(a<b\) we must have \(a\vee p<b\vee p\) by intersection strictness (actual
intersections nontrivial since joins are residual). A representative for \(p\) cannot move between
units. Otherwise for an element sending a letter of \(j\) into \(k\ne j\), choose a second letter
outside \(k\) mapping outside \(k\) (available by sizes with copies). Use a supported full
symmetric label on the one or two source units. Conjugation of the source-letter transposition in
the join gives a supported transposition connecting \(k\) with somewhere else, adjoining which to
the source label is a small strict increase by orbit detection at \(k\), impossible.

\noindent\textbf{Preserving parallelisms.} On a square both parallelisms, and on an auxiliary grid both length-two parallelisms, must be preserved. Otherwise use the parallel partition plus matching reference block with full symmetric wreath action and only \(Y\) elsewhere as \(a\), and the independent full symmetric action on that union as \(b\). Again disruption generates \(b\) in the join with \(p\); strictness is detected by meeting with the other indicated parallel label of equal length using the \emph{same} reference block, then testing whether the reference unit connects to the grid (it is the only common block of the two partitions).

\noindent\textbf{Uniform auxiliary flips.} Thus on an auxiliary unit \(h_0\times2\times2\) a representative acts by component permutations on \(h_0\), and separate bit flips possibly per component. They must actually be uniform. If for instance the first flip varies, take \(a\) acting on the (component, second bit) positions by a full symmetric group and separately by constant first flip, with only \(Y\) elsewhere. Conjugation of its supported top permutations by the varying element produces (after removing top action) first-flip differences generating all even-sum flips on those positions: differences of the nonconstant flip vector with its translates include weight two and are supported within this unit regardless of actions on other units. Adjoining those to \(a\) strictly increases it. Indeed meet with simultaneous preservation of the two parallelisms: the first flips formerly constant induce with component permutations only even actions on the (component, first bit) lines (second-bit direction, \(h_0\) even), but flipping just one component's first bit now occurs and acts oddly. As before distinguish by adjoining a spare marker swap then intersecting a parity graph label, which contains bottom. This violates strict joins again.

\noindent\textbf{The remaining primary signs.} We conclude \(p\le B\) of the earlier sign-space test, with exact \([0,B]\) since \(p>0\). On any primary slot with active sign take as \(a\) the alternating bases on those lines plus the auxiliary size-\(x\) block, with \emph{symmetric} top, outside only \(Y\). In the join with \(p\), conjugation again adds independent augmentation base signs by discrepancy with the even auxiliary action, producing the strict increase \(b\) tested under \(B\) as before. This final contradiction (or triviality already if no signs) eliminates the proposed residual product.
\end{proof}

\begin{proposition}[Reverse projections and the trivial-base cases]\label{t:weighted-reverse-projections}
In reverse orientation, if \(P>1\), then \(c(P)=\mathcal W\).
If \(P=1\), every nonidentity actual label is subdirect and the site-action
image \(\Lambda\) contains the inner automorphisms of \(S\).
In forward orientation with \(P=1\), either a marker axis filter gives
product success, or the same subdirectness and inner-action conclusions hold.
\end{proposition}
\begin{proof}

For \(P>1\) in reverse orientation, a proper closure would be residual by
\Cref{t:weighted-small-labels} and would contradict
\Cref{t:weighted-residual-obstruction}.  This proves the first assertion.

If \(P=1\), apply the result just obtained in the reversed filter with parameter range
\([0,M_{\{i\}}]\) for a spare marker \(i\) outside any tested small support. The actual base there
is nontrivial with still the same transitive operators, so it projects fully by the reverse
\(P>1\) result and the desired small labels are indeed subdirect. A proper nontrivial product
group would have residual label \(p\), since containing a singleton absorption in parameter order
would put independent nontrivial support inside that small proper subdirect group. The residual
argument rules it out, hence all nonbottom groups are subdirect as claimed. Here \(\Lambda\) must
then contain inners (otherwise its proper nontrivial inner part, or the proper-invariant-subgroup
result for soluble operators in \Cref{sec:affine}, yields a proper nontrivial invariant product).

Back forward, suppose \(P=1\). If some marker \(x_i\) does not project fully, its axis filter
gives product success by the forward argument above (we may take the enlarged actual base as
operators then). Otherwise the extension at each marker \(x_i\) has image containing inners by
atomicity and the same proper-invariant-subgroup argument, now on a diagonal part: a nontrivial
proper subgroup invariant there would transport and multiply to an invariant group strictly
between \(P=1\) and \(K_{x_i}\). Its domain interval via restriction realizes the reverse weighted
test on one less marker. If \(\Lambda\) itself does not contain inners, the kernel of that
extension map is retained by the normal test (supplementation would give image just \(\Lambda\)),
hence this gives an ordinary interval within the image in \({\rm Aut}(S)\), with base \(\Lambda\),
top containing \(S\). Missing inners there must be supplemented, so this is a reverse invariant
interval in one simple site with proper base \(\Lambda\cap S\), impossible by the reverse result
just proved: a nontrivial base cannot project fully here, and for trivial base having all
nonbottom actual groups project fully leaves no interior labels. Thus \(\Lambda\) contains inners
and again all nonbottom labels are subdirect.
\end{proof}

\subsection{Extension domains and the passage to product success}

\begin{theorem}[Product success with absolute stock]\label{t:weighted-product-success}
There are absolute stock and marker-loss bounds with the following
property.  Any ordinary realization, in either orientation, of a nontrivial
weighted image yields a forward realization of a nontrivial image of the
weighted template with at most boundedly many markers deleted, in which
product success holds.  The stock bound is independent of all numerical
unit sizes and of the number of cohorts.  The simple group at a site has
order at most the original top-group order.  A larger prescribed fixed
stock can be retained by increasing the initial absolute stock.
\end{theorem}
\begin{proof}
Take an ordinary realization of minimum top order among both orientations
of the same quotient; its order is at most the input order.  If
\(c(P)=\mathcal W\), the ties of \(P\) are fixed and every actual label
refines them.  On the tie part through a chosen site these refinements are
exactly the invariant equivalences under its part stabilizer in \(H\).
They give an ordinary realization of the opposite interval with top order
at most \(|H|<|G|\), as in the proof of \Cref{f:case-subdirect}.
This contradicts minimality. Thus absent product success we have the artificial-top extension
case in the opposite orientation, with base map containing inners.

\noindent\textbf{The extension datum and its flat labels.}
To apply \Cref{f:lem-lower-grid}, distinguish the new base from the
old operator group.  Write \(H_{\rm op}\) for that operator group,
\(A_0\le H_{\rm op}\) for its site stabilizer, and
\(\alpha:A_0\to\operatorname{Aut}(S)\) for its action on the site.
The remaining case gives \(\alpha(A_0)\ge\operatorname{Inn}(S)\).
Thus the extension datum in that lemma is
\[
 (K,H,B,\psi)=(H_{\rm op},A_0,S,\alpha).
\]
Put \(C_0=\ker\alpha\) and
\(A_{\rm in}=\alpha^{-1}(\operatorname{Inn}(S))\); these are its
base kernel and inner preimage.

Index the points of an affine four-space over \(\mathbf F_3\) by
distinct markers.  For a nonempty flat \(X\) of dimension at most two,
let \(Z_X\) be the extension label given by \(M_X\) in forward
extension order and by \(X_X\) in reverse extension order.  Write
\(D_X\) and \(E_X\) for its domain and kernel.  Restriction identifies
the entire lower ideal of this extension with the ordinary interval
\([A_0,D_X]\).  In parameter order this is the axis ideal
\([0,M_X]\), or the axis filter \([X_X,1]\) with order reversed.
It is a nontrivial weighted image deleting just the markers of \(X\),
so \(D_X>A_0\) and its normal property \((*)\) is available.

Apply that property to \(E_X\trianglelefteq D_X\).  If its kernel
were retained in \(A_0\), the whole preimage of
\(\operatorname{Inn}(S)\) would lie in \(A_0\): an element of that
preimage has the same image as some element of \(A_0\), and their
quotient lies in \(E_X\).  This normal subgroup would lie in the base
core, leaving a soluble quotient by Schreier's theorem, contrary to
the nonabelian monolith in \((*)\).  Hence the kernel supplements
the base.  Each flat extension is pure, with
\[
 D_X=A_0E_X,\qquad E_X\cap A_0=C_0.
\]
Only these flat labels need purity.

\noindent\textbf{Proper and incompatible kernel joins.}
The literal map \(D\mapsto X_D\) preserves Boolean joins and meets,
and footprints preserve its strictness in the image.  Consequently
the \(Z_X\)'s reverse flat containment, and for \(X\cap Y=\{i\}\)
their join in extension order is the point label.  Restricting the
point extension to the generated domain gives
\[
 D_i=\langle D_X,D_Y\rangle
     =A_0\langle E_X,E_Y\rangle,
 \qquad E_i=\langle E_X,E_Y\rangle.
\]
The kernel equality follows because the generated kernels lie in
\(E_i\), have the same intersection \(C_0\) with \(A_0\), and
already generate its domain together with \(A_0\).

Distinct point labels have incompatible join.  For their generated
kernel group \(E\), the equality \(E\cap A_0=C_0\) would define a
common pure extension on \(A_0E\), which is impossible.  Hence
\(E\cap A_0>C_0\).  It is normal in \(A_0\), and its nontrivial
\(\alpha\)-image contains the inner simple subgroup.  Since
\(C_0\le E\), this gives \(A_{\rm in}\le E\), exactly the
incompatible-join condition of the lower-grid lemma.

\noindent\textbf{Subdirect intersections in the point monoliths.}
Apply the projection conclusions to each ordinary point ideal
\([A_0,D_i]\), without a minimality assumption.  If they give product
success, the theorem is already proved on the corresponding bounded
marker restriction.  Otherwise a nontrivial base intersection has
full projection: this follows from the reverse projection result in
reverse orientation and from the absence of forward product success
in forward orientation.  With trivial base intersection, every
nonidentity label is subdirect by
\Cref{t:weighted-reverse-projections}, again allowing its marker-filter
success alternative.

Thus the domains of lines and planes through \(i\) intersect the
point monolith subdirectly modulo the point core.  These intersections
are strictly above the point-base intersection: each such domain is
strictly above \(A_0\), by footprints, and supplementation recovers
the domain from its intersection.  The flat ideals have \((*)\),
the proper kernel joins were just verified, and incompatible point
joins contain \(A_{\rm in}\).  All hypotheses of
\Cref{f:lem-lower-grid} therefore hold, giving the required contradiction.

\noindent\textbf{Stock and order bounds.}
Each test uses boundedly many units: four-block comparisons, packages
of at most ten units, or the fixed affine grid.  The Ramsey reservoir
for the copy-action argument concerns package indices and bounded
permutation patterns, independently of the numbers of letters inside
a package.  Cloning a bounded package therefore needs only an absolute
number of copies from any one cohort, including the two consecutive
copy counts used earlier.  Tests on different cohorts reuse these
stocks; they do not delete the tested units.  The lasting deletions
are the boundedly nested marker restrictions in the extension,
atom-filter and point-ideal steps.  Increasing the initial stock thus
leaves any prescribed fixed number of spare copies.

Only core quotients, subgroups and domain intervals have been used;
the extension domains lie in the preceding operator group.  Hence
every resulting simple site has order at most the initial minimal
top-group order, itself at most the input order.  The reverse
projection step uses \Cref{r:reverse-exclusion}.
\end{proof}

\subsection{A single alternating action for all readings}

\begin{lemma}[Restart in an actual alternating power]\label{t:weighted-restart}
From product success one obtains, after boundedly many further marker
deletions, an actual power \(A_\Omega^s\) with transitive operator action,
base intersection \(H\cap A_\Omega^s\), and product subgroups at every cut.
The degree is high, and at least forty specified surviving marker cuts are
genuine distinct alternating point stabilizers at each site.
All previously product labels below the new top remain products.
\end{lemma}
\begin{proof}
In the forward product success case, take the pure filter absorbing everything except markers (one
can leave just a sufficiently large fixed substock unabsorbed and put many spare markers among the
absorbed units). Its base \(K_{\rm sat}\) inside \(\mathcal W\) is product by intersection of
marker point cuts, nontrivial, and the filter exact pure. Apply the forward marker theorem for,
say, at least 40 distinct points. Thus at \(M_A\) for boundedly many omitted markers we have an
actual high alternating simple \(B_A\) at one site in the ordinary product group \(\Lambda
T_{M_A}\) (\(T_b={\rm proj}(K_b)\)), normal and supplemented by the pure base there. Here
\(\Lambda\cap S\le T_{\rm sat}\) as in that theorem; thus \(B_A\le T_{M_A}\), and multiplying its
transported factors gives a normal group at actual \(HK_{M_A}\), not contained in the pure base.
In the whole ideal below \(M_A\), removing these boundedly many markers, (\(*\)) shows that
already \(H\) supplements this normal group.  Explicitly, normality of
the local \(B_A\) under \(\Lambda\) makes its transport independent of
the chosen element of \(H\).  The resulting actual subgroup
\(N=\prod B_A\) lies in the old power, is normalized by
\(K_{M_A}\) and by \(H\), and is not contained in
\(HK_{\rm sat}\), hence not in \(H\).  The whole-ideal normal test
gives \(HN=HK_{M_A}\); the original transitive action of \(H\) on
the sites permutes these very factors.
Restart therefore with its power of alternating
groups \(B_A^s\), transitive operator action, and base intersection \(H\cap B_A^s\); labels that
were products below \(M_A\) remain product by intersection. In particular all relative cuts are
product (intersect original cuts with \(M_A\)). The specified genuine point stabilizers still hold
by intersection with \(B_A\). Below use notation on the remaining template (removing \(A\)) and
this new power.

If success was obtained by an earlier axis-filter shift deleting a marker, take \(H\) to be the
ordinary base of the resulting whole weighted template (thus enlarged at that prior shift). It can
remain the operator group in testing the pure filter \([K_{\rm sat},\mathcal W]^H\); that filter's
\(M_A\) is exactly the lift \(M_A\) in the whole template. In the transport just used, \(K_{M_A}\)
itself is a product by the marker-cut intersections, \(\Lambda\) normalizes the local \(B_A\), and
\(T_{M_A}\) normalizes it as well. Its power is not in the pure saturated base since \(M_A\)
strictly increases saturation and the local simple supplements it on ordinary product labels.
After supplementation by \(H\), intersection with this new power at each label is just the
intersection of the old \(K_b\) with it (the old power contained the new one). For \(M_{A\cup i}\)
from the pure theorem, local intersection with \(B_A\) can be taken on the ordinary site label as
well, since \(\Lambda\cap S\le T_{\rm sat}\); the almost simple quotient's centralizer kernel was
in that ordinary base. Thus it gives exactly the alternating point stabilizer inside the new site.
This step uses the forward marker result and its high-type reductions in
\Cref{sec:forward-completion}.
\end{proof}

\begin{convention}[Ordinary site groups and groups in the power]\label{t:weighted-site-convention}
Write \(B_A=A_\Omega\); locally let \(\Theta\) be the site action image, in \(S_\Omega\) since the
degree is high. The projection of the new base intersection lies in \(\Theta\), and all cut site
groups are nontrivial proper normalized by \(\Theta\). They contain \(\Theta\cap A_\Omega\):
adjoining it there could not reach \(A_\Omega\) by simplicity, hence could not increase by the
product coatom property. We may thus look also at the ordinary cut groups multiplying by
\(\Theta\), inside \(\Theta A_\Omega\). This respects intersections (the cut site groups all
contain \(\Theta\cap A_\Omega\)); covers between such products remain covers as ordinary site
groups since any strict intermediate, intersected with \(A_\Omega\), would transport to a strict
intermediate product.

In particular, throughout the readings \(K_b\) on this restarted template denotes an actual
intersection inside the normal power (before multiplying by any ordinary base). We use the
ordinary cut groups with \(\Theta\) mainly to test permutation actions; when using a projection or
a graph in the power keep the parity restrictions in the alternating groups. A proposed
intermediate \emph{product} obtained sitewise is a label if it contains the base intersection and
is invariant (transport to the other sites is implicit); arbitrary intermediate labels need not
themselves be products.
\end{convention}

\begin{lemma}[All cuts are full nonhalf subset stabilizers]\label{t:weighted-cuts}
In the alternating realization of \Cref{t:weighted-restart}, every cut
coatom has site group the alternating intersection of a full subset
stabilizer on \(\Omega\).  The two parts of this subset cut have unequal
sizes.  The ordinary site groups obtained by multiplication by \(\Theta\)
respect all cut intersections and the indicated covers.
\end{lemma}
\begin{proof}
Every cut coatom \(C_D\) (remaining indices) must give a full nonhalf subset stabilizer on
\(\Omega\). Indeed use at least eight genuine distinct point markers all on one side of the formal
cut, and up to six of them at a time as a prefix \(P_1\). The cut group intersected with this full
pointwise stabilizer (axis \(M_{P_1}\), by meeting the individual full point stabilizers) is
maximal proper there including in ordinary site notation with \(\Theta\). In labels it is a
relative cut coatom after deleting the prefix, and these are products; any ordinary intermediate
would transport to a product. Also the intersection with a single such point stabilizer is covered
by the cut itself, literally by the independent full groups on the two pieces in the template
(strictness as well by footprints or orbit detection, and use the cover transfer between the
indicated products as ordinary site groups). Thus if the ordinary cut group were transitive, it
would be primitive by point stabilizer maximality. Here the tested genuine points are fixed
individually by \(\Theta\) (it normalizes their respective alternating point stabilizers), so
prefix stabilizers indeed cut down inside a pointwise ambient alternating or symmetric group.
Every tested prefix stabilizer in it would have to be transitive on the remaining letters:
otherwise by its maximality in the full pointwise ambient group it would be a full subset
stabilizer there and contain a three-cycle fixing the prefix points, impossible in the original
assumed transitive primitive cut group by Jordan's Theorem (its intersection with \(A_\Omega\) is
proper). Indeed on deleting the prefix the residual cut is still proper with a nonempty unabsorbed
side by our spare markers; for a single chosen marker the upward cover back to the cut is just
restoring one singleton to a full symmetric action on its piece. Strictness if that piece is
unabsorbed follows by absorption on adjoining that marker. Taking successive nested prefixes gives
at least six-fold transitivity, impossible here by the classification of multiply transitive
groups. Thus the cut is intransitive and by maximality a full subset stabilizer (its alternating
intersection is exactly the product site group); a half subset would extend properly to the
transitive two-block stabilizer instead.
The permutation-group inputs in this argument are Jordan's three-cycle
Theorem and the classification of multiply transitive groups, in the
high-degree range, as recorded in \Cref{sec:affine}; see also
\citet[Theorem~1.1 and Section~2]{Jones2014} for the prime-cycle
statement and the exclusion of proper six-transitive groups.
\end{proof}

\begin{proposition}[Common alignment of all cuts]\label{t:weighted-alignment}
There is a single bijection from the formal units together with the
foundation position to the \(\Theta\)-orbits on \(\Omega\), under which
every cut and every intersection of cuts is the corresponding stabilizer
of unions of those orbits.  Write \(u_j>0\) for the size of the cell of unit
\(j\).  The specified genuine point markers have size one.
\end{proposition}
\begin{proof}
\noindent\textbf{Every orbit bipartition is a cut.}
The preceding lemma assigns a bipartition into unions of
\(\Theta\)-orbits to every cut.  Conversely, take any nontrivial such
bipartition.  Its individual-part stabilizer in \(A_\Omega\)
transports to an invariant product containing the base.  A coatom
above this product is a subset stabilizer by
\Cref{t:weighted-cuts}.  It must preserve the same bipartition, since
the original stabilizer together with \(\Theta\) is transitive on
each of its two parts.  Thus the product was already a coatom.
This also excludes a half-sized bipartition.  We have obtained a
bijection between all formal cuts and all bipartitions of the orbit list.

\noindent\textbf{Intersections recover the partition poset.}
At least one \(\Theta\)-orbit is nonsingleton.  Otherwise an even
degree admits a half-sized bipartition.  In odd degree, \(\Theta=1\)
and the base intersection is trivial; a full cycle then transports
to a proper transitive product, whereas every coatom above it would
be an intransitive subset stabilizer, also impossible.

For any partition of the orbit list, intersect the cuts separating
its parts.  At a site this gives the full alternating intersection
of the independent symmetric groups on their unions.  Multiplying
by \(\Theta\) leaves exactly these unions as its orbits.  The only
possible parity exception would be a part of size two with every
other part a singleton: its transposition has no outside parity
correction.  By the nonsingleton-orbit conclusion, that two-point
part is already one \(\Theta\)-orbit, so multiplication by
\(\Theta\) restores transitivity there too.  Hence distinct
partitions give distinct intersection labels, with their correct
containment order.

On the template side, intersections of cuts likewise form the
partition lattice on the units and the foundation position.  A
strict change of partition changes either the absorbed set or a
residual orbit, so footprint and orbit detection preserve strictness
in the image.  The cut-intersection poset therefore identifies these
two partition lattices.

\noindent\textbf{Recovering individual cells.}
A partition lattice on sufficiently many positions has only
position permutations as automorphisms: its merging atoms represent
pairs, their joins detect intersecting pairs, and the large incident
stars recover the individual positions.  The poset identification
thus gives one bijection of formal positions with actual orbit cells,
compatible with every cut and every intersection.  Denote the cell
of unit \(j\) by \(\Omega_j\), and put \(u_j=|\Omega_j|>0\).
Transport these cells to the other simple sites by the operator group.
This is well-defined because the site stabilizer \(\Theta\) preserves
each of its orbit cells; invariance of the cut labels gives the same
formal alignment at every site.
For a specified genuine point marker, the formal singleton cut is
an actual one-point stabilizer.  Its one-cell side must consequently
correspond to that single point, so this marker has size one.
\end{proof}

\subsection{Equality and square readings in that alignment}

The following two readings use the same alignment of \Cref{t:weighted-alignment}.

\begin{proposition}[Equal totals and equal copies]\label{t:weighted-equal-totals}
Let \(D,E\) be disjoint bounded collections of nonfoundation units with
equal formal total size, with ample spare markers outside them.  If at least
one of \(D,E\) contains more than one unit, then
\[
 \sum_{j\in D}u_j=\sum_{j\in E}u_j.
\]
In a sufficiently numerous cohort, all copies have the same actual cell
size.  In particular, all surviving markers have size one.
\end{proposition}
\begin{proof}
Put \(J=C_D\wedge C_E\) and \(F'=C_{D\cup E}\).  Equal formal
totals supply a two-block wreath label \(W\) between them.  It
survives strictly above \(J\): joining an absorption on one side
detects that \(W\) connects the two sides.  To distinguish \(W\)
from \(F'\), intersect with a cut fixing a proper nonempty
subcollection of the multi-unit side.  This prevents the wreath
exchange, while the upper label still mixes letters across the two
sides.  Thus \(J<W<F'\) in the image.

Suppose the actual totals differ.  At one site write \(J_o\) and
\(F_o\) for the projections of the product groups at \(J,F'\).
Projection onto the combined \(D,E\)-cells maps \(F_o\) onto their
full symmetric group: spare units provide outside parity correction.
Its kernel lies in \(J_o\), whose image is the stabilizer of two
unequal nonempty parts.  That stabilizer is maximal and nonnormal,
as is seen by conjugating its supported transpositions.  Therefore
\(J_o<F_o\) is a cover and \(J_o\) is not normal in \(F_o\).

The actual group \(K_W\) need not be a product.  Its coordinate
closure is nevertheless a product label between \(J\) and \(F'\),
so the site cover forces that closure to be \(K_{F'}\).  Now
\(K_W\) contains the independently supported copies of \(J_o\)
and projects onto \(F_o\) at every site.  Conjugating those copies
inside \(K_W\) therefore puts their independent normal closures
in \(F_o\) inside \(K_W\).  Nonnormality and the cover give
\(\langle J_o^{F_o}\rangle=F_o\).  Hence
\(K_{F'}\le K_W\), contradicting \(W<F'\).  The actual totals
must agree.

Apply this conclusion to disjoint pairs of copies in one cohort.
Comparing two pairs which differ at one position against a common
disjoint pair equalizes the individual cell sizes.  This works for
the marker cohort as well; its already identified genuine points
make every surviving marker size one.
\end{proof}

\begin{theorem}[Reading a primary square]\label{t:weighted-square}
For a primary square \(O\) and its full side reference \(R\), let
\(r=u_R\) in the alignment of \Cref{t:weighted-alignment}.  If \(r\) is even
and \(r\ge8\), then \(u_O=r^2\).
\end{theorem}
\begin{proof}
We first construct the two actual parallel systems, then determine their
common incidence cells and the number of blocks.

\noindent\textbf{The two parallel systems.}
Put
\[
 M=M_{\{O,R\}},\qquad F'=C_{\{O,R\}},\qquad
 J=C_{\{O\}}\wedge C_{\{R\}}.
\]
For each of the two formal parallel directions, let \(V_j\) be
the full wreath label whose blocks are the parallel lines and the
reference unit \(R\), with everything outside absorbed.  Each is
maximal below \(F'\): adjoining a permutation disrupting its blocks
and conjugating the supported base transpositions generates the full
symmetric group on their union.  The label \(V_j\) mixes \(R\)
with \(O\), whereas \(V_1\wedge V_2\le J\), since \(R\) is
the only common block of the two formal partitions.  Orbit detection
preserves these distinctions in the image.  In particular both
\(V_j\)'s remain proper and maximal below \(F'\).

First consider the product \(N_R\) of the independently supported
alternating groups on the actual cells \(\Omega_R\) at all sites.
It is normal in the ordinary group at \(C_{\{R\}}\).  Since the
formal reference unit is full, \(C_{\{R\}}=M_{\{R\}}\); the ideal
below it is a weighted image with just that reference unit deleted.
Its normal property \((*)\) applies.  But \(N_R\) lies inside the
separate-cell stabilizer at \(J<C_{\{R\}}\), so it cannot supplement
the base.  It must lie in the base, and hence is available independently
at every site in each \(K_{V_j}\).

Fix a direction \(j\).  The projection of \(K_{V_j}\), together
with \(\Theta\), is transitive on \(\Omega_O\sqcup\Omega_R\):
\(\Theta\) is transitive on each cell, and \(V_j\) mixes them.
Suppose two translates of \(\Omega_R\) overlap without coinciding.
After conjugation one of these translates is \(\Omega_R\) itself.
The component of the overlap relation containing \(\Omega_R\)
then meets \(\Omega_O\) and is \(\Theta\)-invariant.  Cell
transitivity makes it the entire union.  Alternating groups on
intersecting supports of size at least eight generate the alternating
group on their union.  Conjugating the supported factors of \(N_R\)
along the translates therefore gives, inside \(K_{V_j}\), the
independently supported alternating groups on
\(\Omega_O\sqcup\Omega_R\).  Conjugation by a site projection is
realized by an element of \(K_{V_j}\); conjugation by \(\Theta\)
uses operator invariance.  Both preserve single-site support.

Let \(N\) be the product of these larger supported groups.  It is
normal in the ordinary group at \(F'\), since the cell union is
invariant.  In the ordinary interval from the group at \(M\) to
the group at \(F'\), let \(e\) label the lower endpoint multiplied
by \(N\).  We have
\[
 M\le e\le V_j,\qquad e\not\le J.
\]
For the other direction \(k\), set \(a=J\wedge V_k\) and
\(b=V_k\).  Then \(a<b\) and \(e\wedge b\le a\).

The remaining join calculation takes place in the literal template.
Lift \(e\) between \(M\) and \(V_j\).  Because its image is not
below \(J\), this lift moves the block \(R\) to some \(j\)-line.
The literal label \(a\) contains the full supported symmetric
groups on \(R\) and on every \(k\)-line.  Conjugating its group on
\(R\) supplies the full symmetric group on that \(j\)-line.
Its crossings with the \(k\)-lines generate the full symmetric
group on \(O\), and the movement of \(R\) then generates the full
group on \(O\cup R\).  Hence \(e\vee a=F'\).  This is the
forbidden normal-multiple configuration, a contradiction.

Thus the translates of \(\Omega_R\) are disjoint or equal and
form an equivariant block partition at every site.  Its full product
stabilizer is proper, contains \(K_{V_j}\), and lies below \(K_{F'}\).
Maximality forces equality with \(K_{V_j}\).  Removing its block
\(\Omega_R\) gives a partition of \(\Omega_O\) into blocks of
size \(r\).  Doing this for both directions yields two systems,
each with \(m'=u_O/r\) blocks.

\noindent\textbf{The sign graph across all simple sites.}
Set \(a_j=V_j\wedge J\) and \(b_0=a_1\wedge a_2\).  These are full product
subgroups at the simultaneous structures.  Let \(L\) be the projection of
\(K_{b_0}\) at a site onto the actual \(O\)-cell alone.  Outside parity
correction makes this the full common stabilizer of its two systems.
We show that \(L\le A_{u_O}\).

Choose two singleton markers outside \(O,R\).  From \(b_0\) obtain \(b^-\)
by fixing the markers individually and \(b^+\) by fixing their pair setwise;
from \(J\) obtain \(j^+\) by fixing their pair setwise.  Each restriction is
an intersection with the corresponding cut labels.  Formally the group on
\(O\) at \(b_0\) is \(S_x\times S_x\) in its grid action, and every such
permutation is even because \(x\) is even.  The literal graph label \(g\)
inside \(j^+\), which couples the marker swap to the sign on \(O\), therefore
satisfies
\begin{equation}\label{t:weighted-parity-equalities}
 g\vee b^+=j^+,\qquad g\wedge b^+=b^-.
\end{equation}
It does not fix the markers individually: this survives the lattice map
by orbit detection.  This is a valid literal label before taking the
image, since the bottom action on \(O\) is even; actions on the other parts
remain independent.

Now use the actual subgroups in \(A_\Omega^s\).  Project to the \(O\)-cell
and the marker pair at every site, correcting parity on \(R\).  The three
product subgroups \(K_{b^-},K_{b^+},K_{j^+}\) have a full common kernel,
say \(Q_0\), and their quotients are respectively
\[
 L^s\times0,\qquad L^s\times C_2^s,
 \qquad S_{u_O}^s\times C_2^s.
\]
The restrictions distinguishing them concern only \(O\) and the marker
pair, so the same full kernel occurs in all three.  By the meet equality,
\(Q_0\le K_{b^-}\le K_g\).  Put \(Q=K_g/Q_0\).  The same equality gives
\[
 Q\cap(L^s\times C_2^s)=L^s\times0.
\]
In particular, the projection of \(Q\) to \(S_{u_O}^s\) is injective and its
image contains \(L^s\).  The join equality in
\Cref{t:weighted-parity-equalities}, taken with
\(L^s\times C_2^s\), forces that image to be all of \(S_{u_O}^s\).
Indeed, if the image is \(P\), then \(L^s\le P\), so the projected
join is \(\langle P,L^s\rangle=P\); it is the entire product,
not merely a subgroup surjecting to each site separately.
Consequently \(Q\) is the graph of a homomorphism
\[
 f:S_{u_O}^s\longrightarrow C_2^s,
 \qquad f(L^s)=0.
\]
No assumption is made that \(f\) acts separately at the different sites.
It is nonzero: if it vanished, \(K_g\) would lie in the product restriction
of \(K_{j^+}\) fixing the two actual points individually, the intersection
with their individual cut labels.  The literal graph connects the two
formal singleton units whereas that intersection does not, with the same
outside absorption, contradicting orbit detection.

Every homomorphism to an abelian group factors through the sign
abelianization of \(S_{u_O}^s\).  If \(L\) contained an odd permutation,
then the signs of \(L^s\) would fill \(C_2^s\), and \(f(L^s)=0\) would force
\(f=0\).  Hence \(L\) is even.  Transport identifies the actual copies of
\(L\) without changing parity.

The common incidence cells of the two systems on \(O\) now have at most one
letter each: otherwise a transposition in one such cell would belong to
\(L\).  Thus the letters of the \(O\)-cell are edges of a simple
\(r\)-regular bipartite graph, with \(m'\) vertices in each part.  The group
\(\Theta\) is transitive on edges, because it is transitive on the
\(O\)-cell.  In particular, \(m'\ge r\).

\noindent\textbf{The row chain and the column signs.}
The literal interval \([b_0,a_1]\) has exactly four elements.  On
the formal \(O\)-grid its lower endpoint contains the diagonal
base \(\Delta S_x\) and all pure row permutations \(S_x\), while
the upper endpoint is \(S_x\wr S_x\).  An intermediate is therefore
determined by its base kernel.  Intersecting this kernel with
\(A_x^x\) gives a subdirect group containing \(\Delta A_x\).
Its tie partition is invariant under all row permutations, so either
all rows are tied or all are independent.  In the first case the
intersection is \(\Delta A_x\), whose normalizer in the base is
\(\Delta S_x\), since \(A_x\) has trivial centralizer in \(S_x\).
In the second case the kernel contains \(A_x^x\).  Its quotient is
an invariant sign space containing the constant vector.  Taking
differences under transpositions shows that the only possibilities
are the constant line, augmentation, and the full sign space.
Thus the four base kernels are
\[
 \Delta S_x
 <A_x^x\Delta S_x
 <(S_x^x)_{\rm even}
 <S_x^x.
\]
Here \(x\) is even, so the constant line lies in augmentation.
The actions outside \(O\), including on \(R\), are independent
and full at both endpoints.  Passing to the lattice image leaves
at most these four elements.

In the actual realization consider the following chain on the
\(O\)-cell, where even refers to total base parity:
\[
 L<L A_r^{m'}<L(S_r^{m'})_{\rm even}<L S_r^{m'}\le S_r\wr S_{m'}.
\]
All groups in this chain lift with the full common kernel; parity can be
corrected freely outside the \(O\)-cell.  They contain the base intersection
and are \(\Theta\)-invariant, because \(\Theta\) preserves both systems and
the displayed base layers are canonical within the row system.  Thus they
are actual product labels in \([b_0,a_1]\).

The first increase is strict: a nontrivial three-cycle supported on one row
cannot preserve the column partition, since the incidences are simple and
every column has degree \(r>1\).  For the next increase, examine the kernel
of the row action of \(L\).  It can only permute columns with the same
neighborhood, called twins.  A permutation inside one twin class has the
same parity contribution on every row of that neighborhood, and contributes
zero on other rows.  Its row-sign image therefore has rank at most the
number of twin classes.  Edge transitivity implies column transitivity,
so all twin classes have the same size \(t\).  If there are no twins the
rank is zero; otherwise \(t\ge2\), and it is at most
\[
 \frac{m'}t\le\frac{m'}2<m'-1.
\]
Intersecting the displayed groups with the row base and comparing their
sign images modulo \(A_r^{m'}\) proves that the augmentation increase is
strict.  The final increase is strict as well: a row permutation has sign
\((\operatorname{sgn}\sigma)^r=1\) on letters because \(r\) is even; since
\(L\) is even, \(L(S_r^{m'})_{\rm even}\) is even, whereas
\(LS_r^{m'}\) contains an odd row-base element.

There are now four distinct labels in an interval which is an image of a
four-element chain.  They exhaust it, so
\(LS_r^{m'}=S_r\wr S_{m'}\).  In particular the row action of \(L\) is
\(S_{m'}\).  The orbit of any column neighborhood under this row action
contains every \(r\)-subset of the rows, whereas there are only \(m'\)
columns.  Hence
\[
 \binom{m'}r\le m'.
\]
As \(m'\ge r\ge8\), this forces \(r=m'\) or \(r=m'-1\).
In the latter case the missing incidences form a perfect matching:
every row and every column misses exactly one partner.  Label each
column by the row it misses.  The letters then identify with ordered
pairs \((i,j)\) of distinct indices, and every element of \(L\)
acts by the same permutation on both indices.  The full row action
\(S_{m'}\) supplies an element inducing a transposition \((a\ b)\).
On letters it exchanges \((a,b)\) with \((b,a)\), and for each
other index \(c\) it exchanges both \((a,c)\) with \((b,c)\) and
\((c,a)\) with \((c,b)\).  This is a product of
\(1+2(m'-2)=2m'-3\) transpositions, hence odd, contrary to
\(L\le A_{u_O}\).  Thus \(m'=r\), and the \(O\)-cell has
\(m'r=r^2\) letters.
\end{proof}

\Cref{fig:weighted-square} illustrates the incidence conclusion just proved.

\begin{theorem}[Weighted realization and simultaneous numerical readings]\label{t:weighted-realization}
Fix any bounded support size for the comparisons and any desired fixed
number of spare copies.  There are absolute initial cohort-stock and
marker-loss bounds, independent of the numerical unit sizes and the number
of cohorts, with the following property.  The bounds may depend on the
fixed support and spare-copy requirements.  An ordinary realization of a
nontrivial weighted image, in either orientation and with top order \(M\),
gives a forward realization after at most that many marker deletions, with
all cuts product in an actual alternating power.  A single common cell
alignment gives positive sizes \(u_j\le M\), equal on each cohort and one
on all surviving markers.  In this same alignment the equal-total reading
of \Cref{t:weighted-equal-totals} and the square reading of
\Cref{t:weighted-square} hold.

If the input is a finite algebra of size \(m\) representing a nontrivial
image, the same conclusion holds with \(M=60^m m!\), for a possibly
different nontrivial image of the original template.  The full template
itself is finitely representable.
\end{theorem}
\begin{proof}
Apply \Cref{t:weighted-algebra,t:weighted-product-success,t:weighted-restart},
then \Cref{t:weighted-alignment,t:weighted-equal-totals,t:weighted-square}.
All sizes read here are bounded by the order bound already given, since the high alternating
occurs inside an initial site. Stock bounds, including losses to \(A\) and marker requests, remain
independent of the numerical sizes read.
In particular \(u_j\le |\Omega|\le |A_\Omega|\le M\).  All comparisons are
made in this one realization; none requires a new choice of simple sites
or a new alignment.  The original stock can include the fixed losses and
the fixed spare copies required at the end.
\end{proof}

\begin{figure}[htbp]
\centering
\begin{minipage}{0.77\linewidth}
\centering
\begin{tikzpicture}[x=0.5cm,y=0.5cm,
  every node/.style={font=\small},
  letter/.style={circle,fill=black,inner sep=1.05pt}]
  \fill[black!9] (0,5) rectangle (8,6);
  \fill[black!9] (2,0) rectangle (3,8);
  \draw[step=1,black!25,thin] (0,0) grid (8,8);
  \draw (0,0) rectangle (8,8);
  \draw[thick] (0,5) rectangle (8,6);
  \draw[thick,densely dashed] (2,0) rectangle (3,8);
  \foreach \i in {0,...,7}{
    \foreach \j in {0,...,7}{
      \node[letter] at (\i+0.5,\j+0.5) {};
    }
  }
  \node[anchor=east,align=right] at (-0.6,5.5)
    {one row\\\(8\) letters};
  \draw[-{Stealth[length=3pt]}] (-0.5,5.5)--(-0.08,5.5);
  \node[align=center] at (2.5,9.15)
    {one column\\\(8\) letters};
  \draw[-{Stealth[length=3pt]}] (2.5,8.55)--(2.5,8.08);
  \node[align=center] at (4,-0.85)
    {actual \(O\)-cell\\\(u_O=8^2\)};

  \draw (10.4,0) rectangle (11.4,8);
  \foreach \j in {1,...,7}{
    \draw[black!25,thin] (10.4,\j)--(11.4,\j);
  }
  \foreach \j in {0,...,7}{
    \node[letter] at (10.9,\j+0.5) {};
  }
  \node[align=center] at (10.9,9.15)
    {reference\\cell \(R\)};
  \node[align=center] at (10.9,-0.85)
    {\(u_R=8\)};
\end{tikzpicture}
\caption{The established conclusion of \Cref{t:weighted-square},
illustrated for actual size \(r=8\) at one simple site.  In the common
cell alignment, the \(O\)-cell has \(r\) row blocks and \(r\) column
blocks, each with \(r\) letters; each intersection contains one letter
(a dot).  The separate reference cell \(R\) contains \(r\) letters.}
\label{fig:weighted-square}
\end{minipage}
\end{figure}

\FloatBarrier
